\chapter{Renacimiento y revelaciones}\label{sec:renacimiento-y-revelaciones}

Muerte clínica. Una expresión humana que se refiere a la ausencia de respiración y de latidos cardíacos. Lo que cualquier especie normal llamaría simplemente muerte.

El médico del ala médica del buque insignia me dijo que estuve sin funciones vitales durante poco más de dos minutos. Esta afirmación parecía imposible, dado que estaba muy vivo. Cuando lo presioné para que me explicara cómo era posible, básicamente dijo que habían devuelto mi cadáver a la vida, aunque disfrazó el concepto con jerga científica. Creo que el término era «compresiones torácicas». Si esa era la idea terrana de un chiste, no tenía ninguna gracia.

Esa herida debería haber sido mortal según todos los estándares médicos conocidos. Al despertarme en una habitación surrealista, completamente blanca, dopado con cualquier cóctel de fármacos que estuviera en esas vías intravenosas, mi primer pensamiento fue que estaba muerto. Todavía no estaba completamente convencido de que esto no fuera el más allá, o algún truco de mi mente agonizante.

 —¿En qué piensa? —era el comandante Rykov, sentado en el taburete junto a mí, vestido de civil. El buque insignia había hecho escala en la Luna para realizar tareas de mantenimiento después de nuestra pequeña aventura, y yo había aceptado su oferta de ir al bar local—. Se quedó mirando fijo el mostrador.

Intenté salir de mis pensamientos y suspiré. —Creo que necesito un trago.

Se rió entre dientes y le hizo una seña a la mesera. —Ese es el espíritu. Dos tragos de vodka.

La mesera asintió y se alejó arrastrando los pies para atender el pedido. Sentí que algunas personas me miraban y me di cuenta, con cierta incomodidad, de que yo era el único no humano en el establecimiento. Miré las pantallas de televisión en la pared, tratando de distraerme. El sonido estaba silenciado, pero podía ver los subtítulos.

\textit{...las peticiones de dimisión de la presidenta Ula se han intensificado tras filtrarse a la prensa memorandos internos que denigran a varios representantes de la Federación. El origen de esta información no está claro, pero ella ha acusado directamente a la Unión Terrana.}

 —¿Qué dijo exactamente? —pregunté, señalando la pantalla.

—Bueno, para empezar, llamó al embajador Aroktar «cobarde, estúpido y posiblemente descerebrado». Y luego procedió a decir que difícilmente consideraría a su especie como sintiente. Puede ver por qué eso podría ser ofensivo.

—¿Podría ser? Yo diría que sus días como presidenta están contados —miré fijamente al comandante y me crucé de brazos—. Me alegro de que los amigos de su hermano la hayan delatado.

Rykov se puso rígido y la alarma brilló en sus ojos. —No sé de qué...

—Le detendré ahí mismo. Sí lo sabe —gruñí—. ¿Cree que soy estúpido? Todos pensábamos que los hoda'al eran simplemente paranoicos, pero tenían razón cuando dijeron que sus diplomáticos eran espías, ¿no?

—Entre nosotros, sí. Pero es mejor no hablar de esto. La Agencia no tiene reglas y tiene oídos en todas partes —el humano había bajado la voz a poco más que un susurro. Nunca lo había visto tan incómodo, lo que también me puso nervioso—. Tenemos que hablar de todos modos. Pavel... quiere que le pida que guarde silencio sobre cómo murió Cazil. Le dije que podíamos confiar en usted.

En sus palabras había una pregunta implícita: si esa confianza estaba justificada. Su tono era casi suplicante y sus cejas estaban fruncidas por la preocupación. Sospeché que esa preocupación era por mi bienestar, no por el suyo. Si esos espías se metían en medio de un enfrentamiento y asesinaban a sangre fría a un embajador planetario, no tenía ninguna duda de que se asegurarían de que yo guardara silencio, de una forma u otra.

—Puede hacerlo. Mis labios están sellados —dije—. Pero la próxima vez, dímelo sin rodeos. No me gustan esas tonterías del «Departamento de Estado».

—Lo sé, general, y lo lamento. Tenía que mantener su tapadera, pero nunca debí haberlo arrastrado a ese lío en la embajada.

—No me lo hubiera perdido por nada del mundo. Ojalá hubiera podido ver la cara que pusieron cuando apareció el buque insignia.

La mesera regresó y colocó dos vasos diminutos de líquido transparente frente a nosotros. Entrecerré los ojos e inspeccioné las bebidas con sospecha. Era una porción bastante pequeña para una bebida alcohólica, pero ¿quizás los humanos eran más susceptibles a los efectos del alcohol? Sin embargo, después de percibir el aroma de la bebida, retrocedí asombrado. Parecía que era todo lo contrario.

 —¿Está diluida? Huele a alcohol puro.

—Dijo que necesitaba un trago. Pensé que necesitaba algo fuerte, no solo una cerveza —el comandante se bebió todo el vaso de un solo trago y luego me sonrió—. Ahora es su turno.

Dudé, pero tragué el contenido como hizo el humano. La potencia de la bebida me hizo toser; mis ojos se llenaron de lágrimas por su intensidad. Me llevé la mano a la garganta mientras sentía el vodka quemando un rastro a su paso. Lo único que evitó que me cayera de la silla fue el comandante Rykov, que parecía intacto cuando me sujetó.

 —¿Qué les pasa a ustedes? —balbuceé.

Parecía que mi lucha había captado la atención de los clientes del bar, y muchos se reían a mi costa. Al ver la alegría genuina en los ojos del comandante, no pude evitar esbozar una sonrisa.

Sacudí la cabeza, intentando despejar mi mente de esa sensación difusa. —Estoy bastante seguro de que eso se consideraría veneno en la Federación.

—La Federación es aburrida. ¿Otro? —Rykov se encogió de hombros.

—Estoy bien.

—Usted se lo pierde.

Me quité los brazos del humano de encima e intenté ponerme de pie. El mareo no era precisamente agradable, así que caminar me pareció una buena idea. Solo di unos pasos antes de tambalearme. Intenté apoyarme en el mostrador pero fue en vano; lo único que logré fue cortarme la mano con el borde.

Hice una mueca ante el escozor y me miré la palma. Sentí una profunda impresión al ver que la herida se encogía ante mis ojos y se cerraba como por arte de magia. El dolor disminuyó hasta convertirse en un suave hormigueo, mientras la zona se calentaba de forma incómoda.

Eso no era normal. ¿Quizás fue el vodka?

No, el alcohol no funcionaba así. Si tan solo pudiera pensar con claridad ahora mismo...

Me quedé mirando al comandante Rykov, esperando que tuviera alguna explicación. No habría creído algo así si no lo hubiera visto con mis propios ojos. Sentí un escalofrío al notar que el humano parecía más alarmado que sorprendido. ¿Qué habían hecho exactamente para salvarme la vida? Algo me decía que no me habían contado toda la verdad. Una vez más.

 —¿Qué hizo? —escupí.

—Mire, fue mi culpa que le dispararan. Tenía que arreglarlo. Tenía que hacerlo —murmuró—. Le inyectamos nanotecnología experimental.

Respiré entrecortadamente unas cuantas veces, intentando procesar la información. ¿Los humanos habían convertido los nanitos en armas y luego los habían usado con fines medicinales? Todo el asunto sonaba ilegal y peligroso. Pero, dado que yo habría muerto sin la intervención, ¿realmente tenía derecho a quejarme? No parecía haber hecho ningún daño.

De hecho, una muestra de mi sangre podría ser muy útil para la Confederación Jatari. Sospeché que si las dos especies militares restantes querían seguir siendo relevantes, teníamos que ponernos al día con la tecnología terrestre, o de lo contrario la Federación podría recurrir a la Tierra como su único protector. La ingeniería inversa aceleraría ese proceso por décadas.

En ese momento, la Unión Terrana era un amigo y esperaba que siguiera siendo así. Pero los tiempos y los gobiernos cambian, y estar a merced de los cambios de humor de la humanidad no parecía la mejor decisión de política exterior. Era probable que la Confederación intentara convertirse en el nuevo socio comercial de la Tierra y, después de presenciar la implosión de la República Xanik, sabía que debíamos adoptar un enfoque más cauteloso. No enredarnos demasiado ni depender demasiado de ella.

—General, por favor, hábleme —la voz del comandante era suplicante, casi desesperada—. Le diré todo lo que quiera saber.

Me sacudí el polvo y me levanté con las piernas temblorosas. —Quiero saber qué es lo que hacen con los nanitos.

—Usted es un ejemplo —se rió y me dio una palmadita en la espalda.

Sonreí, pero mi mente estaba en otra parte. Tal vez sería mutuamente beneficioso si negociáramos algún tipo de trato bajo la mesa. Podríamos hacer avanzar la investigación de los humanos, ofrecerles una generosa cantidad de créditos, a cambio de armas y tecnología. Ahora que me di cuenta de lo mucho que les gustaba a los terranos la sutileza, pensé que solo necesitaba lograr que mi gobierno se uniera a ellos en las sombras.

Decidí abordar el tema con el comandante más tarde, cuando estuviera sobrio y estuviéramos en un lugar menos público. Con un poco de suerte, nuestra amistad daría lugar a una asociación duradera entre nuestras especies. Como iguales.